\section{红队结论：复杂框架不能自动产生优势}

本报告最需要防范的是一种舒适感：因为有了全球机构名称、公式、表格、风险术语和流程图，就误以为个人已获得机构优势。事实上，流程只能减少可避免的错误，不能消除估值风险、宏观断裂、政策变化、信息时滞、交易受限、模型误设或行为偏差。任何一个被称为“长期”“量化”“分散”或“另类”的框架，都应该先通过以下反方场景。

\begin{exhibitbox}[Exhibit 13：机构策略翻译的红队场景]
\small
\begin{tabularx}{\exhibitboxwidth}{L{2.8cm}X L{3.85cm}}
\toprule
\textbf{反方场景} & \textbf{为什么会使看似合理的框架失效} & \textbf{本报告要求的反应} \\
\midrule
趋势长期震荡 & 多速度 CTA 在持续趋势缺失的阶段会反复换向，成本和滑点可能主导结果。 & 保存每次换向、成本、无成交与换月；不因近期回撤临时改参数。 \\
共同因子拥挤 & 多个因子、行业或宏观袖套在流动性收缩时可能同时受损。 & 按收入、融资、商品、利率、客户和政策链条穿透，不按代码数量判断分散。 \\
财报口径/预期失真 & 一次性项目、修订、披露延迟或不可审计预期会制造虚假的“惊喜”。 & 将事件降级为背景资料；可得时间和预期来源缺失时不计算分数。 \\
市场无法处置 & 停牌、价格限制、深度下降、保证金上调或交割限制使预设退出不可实现。 & 把未成交作为结果，冻结新增；不按理想价格补算。 \\
机构能力错配 & 借券、融资、场外工具、做市库存、专属数据或控制权是策略关键，个人研究通常缺少这些能力。 & 将该策略标为不可等价复制；名称相似的替代品不能解决该能力错配。 \\
监管/账户变化 & 产品、合约、资格、税务、场所或数据许可可能变化。 & 规则变化即重新核验，旧结论随之失效；研究不设计规避路径。 \\
行为性过拟合 & 研究者只记住成功案例、对失败做解释、频繁调参并把复杂度当成熟。 & 冻结版本、保留失败日志、由非作者审查，必要时永久归零。 \\
\bottomrule
\end{tabularx}
\sourcenote{S01--S15；红队情景用于本报告的反方推演，未对未来概率或收益作出预测。}
\end{exhibitbox}

\section{针对三类纸面规则的反例}

\begin{dashboardbox}[CTA、因子、事件各自最容易犯的错误]
\begin{tabularx}{\textwidth}{L{2.1cm}X X}
\toprule
\textbf{规则} & \textbf{最强反例} & \textbf{红队问题} \\
\midrule
CTA & 一段连续图形看似可交易，但实际合约、换月、保证金、限仓和无成交记录并不完整。 & 若删掉连续价格、只保留实际合约与真实限制，结论是否仍存在？ \\
因子 & 数据库用修订后财务字段和存活股票构建，组合看似稳定却缺少当时可知的基础。 & 若按发布日期、停牌/退市和交易成本重放，分数和结论是否仍成立？ \\
事件 & 事后选择了明显的“利好公告”，而当时缺少可审计预期或可处理的市场状态。 & 如果遮住公告后价格并只给原始文件与当时预期，研究者能否仍写出相同规则？ \\
\bottomrule
\end{tabularx}
\end{dashboardbox}

\section{研究停止的触发情形}

以下情形触发研究停止：核心经济机制无法用公开资料解释；数据权利或时点不可核验；策略依赖尚未证明的市场流动性或账户权限；一个结论只有在排除亏损期、无法成交日或错误事件后才成立；或者红队能提出更简单、证据更强的解释但作者无法反驳。研究者不得为此寻找更宽松的参数；此时应维持 \qzero{}，不等待下一个信号。

\section{最终判断}

个人账户不应被视作小型对冲基金、私募股权基金或加密交易台。全球机构策略对个人的启示在于：面对任何策略时，应更早提出困难问题：回报机制是什么？我是否拥有策略成立所需的权利和能力？数据在当时可得吗？最坏情景下能否处置？反方如何证明我错了？这些问题的答案缺失时，零数量比精致的故事更专业。

\begin{riskbox}[最后的非执行声明]
本报告不推荐任何股票、基金、期货合约、虚拟资产、策略、方向、持有期限或配置比例。所有对象均保持 \researchonly{}、\qzero{} 和 \noaction{}；任何将本文转化为外部订单、账户操作、跨境安排或个人适当性结论的行为均超出本报告范围。
\end{riskbox}
